% GENERATED FILE. Edit canonical lessons, then run npm run generate:books.
% canonical-source-hash: fnv1a64:ec640a02452cf5a2
% canonical-lessons: ZH-C487-hangye, ZH-C487-shiye, ZH-W487-ye-career, ZH-C487-jiuye, ZH-C487-shiye-jobless, ZH-C487-gongchang

\chapter{A Career}
\label{ch:zh-a-career}
\hlchaptermodality{487}

\begin{chapteropening}
\textbf{By the end of this chapter:} \emph{I can talk about lines of work, a career, finding a job, being out of work and factories.}

Complete the last lesson of chapter 487: I can talk about lines of work, a career, finding a job, being out of work and factories.
\end{chapteropening}

\section[hángyè]{hángyè — line of business; industry}

\label{lesson:ZH-C487-hangye}



\begin{quote}
Before the new one: say the Mandarin for to do an internship, then the Mandarin for pay and conditions.
\end{quote}

\subsection*{You'll want to know: hángyè}
\emph{hángyè} — \textquotedblleft{}line of business; industry\textquotedblright{}.

\begin{quote}
\emph{nǐ zài nǎge hángyè gōngzuò} — What line of work are you in?
\end{quote}

Its characters come later; for now it is a word for the ear.

\begin{cousinweb}
\emph{hángyè} is built from \emph{háng} 'a row' and \emph{yè} 'a trade', both inherited from Middle Chinese, \emph{hang} and \emph{ngjaep} (Guangyun, in Baxter's transcription), which glosses the first 'a file of soldiers' and the second 'affairs'.
\end{cousinweb}

\begin{grammarlens}[title={What you've built}]
One more word for saying how it was.
\end{grammarlens}

\subsection*{Your turn}
\begin{itemize}
\raggedright
  \item \emph{Say it:} \emph{hángyè}
  \item \emph{Say it:} \emph{hángyè}, once more
  \item \emph{Say it:} \emph{nǐ zài nǎge hángyè gōngzuò} — What line of work are you in?
  \item \emph{Recall:} say the Mandarin for to do an internship, then the Mandarin for pay and conditions, then say \emph{hángyè} again
\end{itemize}

\begin{tcolorbox}[breakable,colback=teal!4,colframe=teal!35!black,title={Before you move on}]
What does hángyè mean? (\textquotedblleft{}Line of business; industry\textquotedblright{}.) Say it once more.
\end{tcolorbox}
\section[\texorpdfstring{\zh{事业} (shìyè)}{事业 (shìyè)}]{\zh{事业} (shìyè) — career; life's work}

\label{lesson:ZH-C487-shiye}



\begin{quote}
Before the new one: say the Mandarin for pay and conditions, then the Mandarin for a line of business.
\end{quote}

\subsection*{You'll want to know: \zh{事业}}
\textbf{\zh{事业}} — \emph{shìyè} — \textquotedblleft{}career; life's work\textquotedblright{}.

\begin{quote}
\textbf{\zh{他觉得事业最重要}} \emph{tā juéde shìyè zuì zhòngyào} — He thinks his career matters most.
\end{quote}

The new character, in pieces: \zh{业} — five strokes and no piece you have written yet: two verticals side by side, a dot on the left and a short stroke falling left on the right, then a long bar along the foot. The next lesson puts it in your hand.

\begin{cousinweb}
\textbf{\zh{事业}} (\emph{shìyè}) is built from \zh{事} 'a matter' and \zh{业} 'a trade', both inherited from Middle Chinese, \emph{dzriH} and \emph{ngjaep} (Guangyun, in Baxter's transcription), which glosses the first 'to employ' and the second 'affairs' and 'great'.
\end{cousinweb}

\begin{grammarlens}[title={What you've built}]
One more word for saying how it was.
\end{grammarlens}

\subsection*{Your turn}
\begin{itemize}
\raggedright
  \item \emph{Say it:} \emph{shìyè}
  \item \emph{Say it:} \emph{shìyè}, once more
  \item \emph{Say it:} \emph{tā juéde shìyè zuì zhòngyào} — He thinks his career matters most.
  \item \emph{Recall:} say the Mandarin for pay and conditions, then the Mandarin for a line of business, then say \emph{shìyè} again
\end{itemize}

\begin{tcolorbox}[breakable,colback=teal!4,colframe=teal!35!black,title={Before you move on}]
What does \zh{事业} mean? (\textquotedblleft{}Career; life's work\textquotedblright{}.) Say it once more.
\end{tcolorbox}
\section[\texorpdfstring{\zh{业} (yè)}{业 (yè)}]{\zh{业} — two verticals, two short strokes, a long bar}

\label{lesson:ZH-W487-ye-career}



\begin{quote}
The word first, then the character inside it.

\begin{itemize}
\raggedright
  \item \emph{Recall:} say \emph{shìyè} — \textbf{R1}, one lesson back, before it decays
\end{itemize}

\begin{quote}
\zh{事} \zh{作}
\end{quote}

\emph{Read it:} each one aloud, with nothing said first — \emph{shì}, \emph{zuò}
\end{quote}

\subsection*{Script — two verticals, two short strokes, a long bar}
\hlblockfigure{\detokenize{figures/ZH-W487-ye-career-filmstrip.pdf}}{How \zh{业} is written, stroke by stroke}

\begin{quote}
\zh{业}
\end{quote}

\textbf{Five strokes, in three steps.}

\textbf{Two verticals first}, side by side in the middle: the left one, then the right one, which starts a little higher.

\textbf{Two short strokes outside them}: a dot on the left that falls to the right, then a short stroke on the right that falls to the left.

\textbf{Last, a long bar} along the foot, wider than everything above it.

2 + 2 + 1 = five.

\subsection*{Your turn}
\begin{itemize}
\raggedright
  \item \emph{Write it:} \zh{业}, all five, with the model in front of you
  \item \emph{Cover:} write it again from memory, and count the pieces, not the strokes
\end{itemize}

\begin{tcolorbox}[breakable,colback=teal!4,colframe=teal!35!black,title={Before you move on}]
How many strokes does \zh{业} have? (\textbf{Five}.) What comes first? (\textbf{The two verticals}.) And last? (\textbf{The long bar along the foot}.)

Source: \href{https://raw.githubusercontent.com/chanind/hanzi-writer-data/68d10a4b21150cae5e1ebbd223eed289cf32d90c/data/\%E4\%B8\%9A.json}{Hanzi Writer Data \zh{业}, pinned stroke paths and medians}.
\end{tcolorbox}
\section[\texorpdfstring{\zh{就业} (jiùyè)}{就业 (jiùyè)}]{\zh{就业} (jiùyè) — to find employment; to get a job}

\label{lesson:ZH-C487-jiuye}



\begin{quote}
Before the new one: say the Mandarin for a line of business, then the Mandarin for one's life's work.
\end{quote}

\subsection*{You'll want to know: \zh{就业}}
\textbf{\zh{就业}} — \emph{jiùyè} — \textquotedblleft{}to find employment; to get a job\textquotedblright{}.

\begin{quote}
\emph{xiànzài dàxuésheng jiùyè bù róngyì} — It isn't easy for graduates to find jobs these days.
\end{quote}

\begin{cousinweb}
\textbf{\zh{就业}} (\emph{jiùyè}) is built from \zh{就} 'to go to' and \zh{业} 'work', both inherited from Middle Chinese, \emph{dzjuwH} and \emph{ngjaep} (Guangyun, in Baxter's transcription), which glosses the first 'to accomplish' and 'to go towards'.
\end{cousinweb}

\begin{grammarlens}[title={What you've built}]
One more word for saying how it was.
\end{grammarlens}

\subsection*{Your turn}
\begin{itemize}
\raggedright
  \item \emph{Say it:} \emph{jiùyè}
  \item \emph{Say it:} \emph{jiùyè}, once more
  \item \emph{Say it:} \emph{xiànzài dàxuésheng jiùyè bù róngyì} — It isn't easy for graduates to find jobs these days.
  \item \emph{Recall:} say the Mandarin for a line of business, then the Mandarin for one's life's work, then say \emph{jiùyè} again
\end{itemize}

\begin{tcolorbox}[breakable,colback=teal!4,colframe=teal!35!black,title={Before you move on}]
What does \zh{就业} mean? (\textquotedblleft{}To find employment; to get a job\textquotedblright{}.) Say it once more.
\end{tcolorbox}
\section[shīyè]{shīyè — to lose one's job; to be out of work}

\label{lesson:ZH-C487-shiye-jobless}



\begin{quote}
Before the new one: say the Mandarin for one's life's work, then the Mandarin for to find employment.
\end{quote}

\subsection*{You'll want to know: shīyè}
\emph{shīyè} — \textquotedblleft{}to lose one's job; to be out of work\textquotedblright{}.

\begin{quote}
\emph{tā shīyè le bàn nián jīngshen hěn bù hǎo} — He was out of work for six months and was in very low spirits.
\end{quote}

Its characters come later; for now it is a word for the ear.

\begin{cousinweb}
\emph{shīyè} is built from \emph{shī} 'to lose' and \zh{业} 'work', both inherited from Middle Chinese, \emph{syit} and \emph{ngjaep} (Guangyun, in Baxter's transcription), which glosses the first 'to err' and 'to let slip'.
\end{cousinweb}

\begin{grammarlens}[title={What you've built}]
One more word for saying how it was.
\end{grammarlens}

\subsection*{Your turn}
\begin{itemize}
\raggedright
  \item \emph{Say it:} \emph{shīyè}
  \item \emph{Say it:} \emph{shīyè}, once more
  \item \emph{Say it:} \emph{tā shīyè le bàn nián jīngshen hěn bù hǎo} — He was out of work for six months and was in very low spirits.
  \item \emph{Recall:} say the Mandarin for one's life's work, then the Mandarin for to find employment, then say \emph{shīyè} again
\end{itemize}

\begin{tcolorbox}[breakable,colback=teal!4,colframe=teal!35!black,title={Before you move on}]
What does shīyè mean? (\textquotedblleft{}To lose one's job; to be out of work\textquotedblright{}.) Say it once more.
\end{tcolorbox}
\section[gōngchǎng]{gōngchǎng — factory}

\label{lesson:ZH-C487-gongchang}



\begin{quote}
Before the new one: say the Mandarin for to find employment, then the Mandarin for to be out of work.
\end{quote}

\subsection*{You'll want to know: gōngchǎng}
\emph{gōngchǎng} — \textquotedblleft{}factory\textquotedblright{}.

\begin{quote}
\emph{wǒ bàba zài gōngchǎng gōngzuò le èrshí nián} — My dad worked in a factory for twenty years.
\end{quote}

Its characters come later; for now it is a word for the ear.

\begin{cousinweb}
\emph{gōngchǎng} is built from \zh{工} 'work' and \emph{chǎng} 'a shed', both inherited from Middle Chinese, \emph{kuwng} and \emph{tsyhangX} (Guangyun, in Baxter's transcription), which glosses the first 'skilled work' and the second 'a roofed building'.
\end{cousinweb}

\begin{grammarlens}[title={What you've built}]
One more word for saying how it was.
\end{grammarlens}

\subsection*{Your turn}
\begin{itemize}
\raggedright
  \item \emph{Say it:} \emph{gōngchǎng}
  \item \emph{Say it:} \emph{gōngchǎng}, once more
  \item \emph{Say it:} \emph{wǒ bàba zài gōngchǎng gōngzuò le èrshí nián} — My dad worked in a factory for twenty years.
  \item \emph{Recall:} say the Mandarin for to find employment, then the Mandarin for to be out of work, then say \emph{gōngchǎng} again
\end{itemize}

\begin{tcolorbox}[breakable,colback=teal!4,colframe=teal!35!black,title={Before you move on}]
What does gōngchǎng mean? (\textquotedblleft{}Factory\textquotedblright{}.) Say it once more.
\end{tcolorbox}
